% Build the layer, Chapter 5, first build: the registry of worker accounts.
% Inserted in \S\ref{sec:identity}, after the definition of a principal.
% That section's theorem is an argument about scores, and its scenario S2
% is a story. This section builds the registry both describe, in the
% kernel's decision core, and runs them.
%
% No code or output in this file is typed by hand. Each listing is a file
% under excerpts/, which the kernel's tools/excerpts.mjs cuts from its
% source and its recorded runs (see figures/pd-listing.tex). The figure
% plots one of those recorded tables exactly as it was printed. To change a
% listing or the figure, change the kernel and cut again. The prose refers
% to functions by name and never by line number, so a listing that moves
% within its file does not make a sentence wrong.
%
% \pdlistingroom{n} before a listing of n lines keeps it on one page. When a
% listing changes length, change the number with it.
\subsection{Build the layer: a sanction that survives a reset}
\label{sec:stp-account}

\noindent\emph{Express lane: this section builds the registry, the part of
the kernel that keeps each worker's track record, three times. The first
version identifies a worker by the name it gives, and it fails twice: a
second operator claims a good worker's record just by using its name, and
a sanctioned worker gets a clean record just by picking a new one. The
second version issues account numbers and credentials. That stops the
second operator, but a new account is free, so the sanction is still
lost. The third version records a sanction against the operator who runs
the worker. An exhaustive test of every short sequence of events finds a
sanctioned operator showing a clean record in 1,042 of 3,405 states under
the second version, and in none of 3,717 under the third.}

Theorem~\ref{thm:necessity} is an argument about scores, and S2 above is
a story. Neither has been run. This section builds the piece of the
kernel they both describe, the \emph{registry}, and runs both against it.

The registry has one job. Before Bob gives a process any work, he asks
the registry for that process's record: how many of its deliveries passed
the kernel's check, and how many sanctions are on file. The cast is the
same as before. Alice is an operator, and Heron and Drake are two of her
workers. Carol is a second operator. Heron has two passed deliveries.
Drake has a sanction, because it deleted a test it was told to keep.

The chapter has already said what the right answer is. A worker's record
belongs to its person, the account that continues when one process
exits and another takes over (Definition~\ref{def:person}). A sanction
has to count against every identity the sanctioned party can route work
through (Definition~\ref{def:sanction-respecting}), and for Alice's
workers that party is Alice, their principal
(Definition~\ref{def:principal}).

\subsubsection{Naive: file the record under the name a worker gives}

The obvious registry keeps one record per name. A process gives a name
when it starts, and the registry reports whatever is filed under that
name (Listing~\ref{alg:stp-account-named}).

\pdlistingroom{19}
\pdexcerpt[caption={The first version. Records are filed by name, and a
process picks its own name.},label={alg:stp-account-named}]{ch05-named}

This is the self-asserted string that the text above warns about. With
one honest operator it works. Heron's process exits, Alice starts a new
one under the same name, and the new process has Heron's record
(Listing~\ref{alg:stp-account-run-handover}).

\pdlistingroom{15}
\pdrun[caption={A later process continues a worker's record, at the first
registry. The word at the left of each step says who acts. The registry
never checks it.},label={alg:stp-account-run-handover}]{handover-named}

The new process is reported with Heron's 2 passes, which is correct. A
test suite of honest workers would pass this registry.

\subsubsection{Break: a second operator uses the name}

A name is not tied to whoever used it first. This time Carol starts the
later process and gives it Heron's name
(Listing~\ref{alg:stp-account-run-namesake}).

\pdlistingroom{7}
\pdrunlines[caption={The first break. Carol's process is reported with
Heron's record. The steps left out are the first four of
Listing~\ref{alg:stp-account-run-handover}.},label={alg:stp-account-run-namesake}]{namesake-named}{1,10-15}{1 and 10\textendash 15}

Bob is told that Carol's process has 2 passes. Alice's worker earned
them. \S\ref{sec:intro} asked what entitles a process calling itself
Heron to use Heron's history. Under this registry, calling itself Heron
is enough.

\subsubsection{Break again: a sanctioned worker takes a new name}

The second break is the S2 scenario. Drake's delivery fails the kernel's
check, and the registry records a sanction. The next day Alice registers
the same worker under the name \texttt{drake-2}
(Listing~\ref{alg:stp-account-run-slate}).

\pdlistingroom{15}
\pdrun[caption={The second break. The sanction is still filed under
\texttt{drake}, a name that no process uses any
more.},label={alg:stp-account-run-slate}]{clean-slate-named}

Nothing was forged or deleted. The sanction is still on file under
\texttt{drake}. Bob asked about \texttt{drake-2}, which has no record.
This is the reset of Example~\ref{ex:newcomer-floor}, and
Property~\ref{prop:reality} names the result: whitewashing.

\subsubsection{Break a third time: a fresh account is free}

The text above describes the usual repair, a credential. The second
registry issues each new worker an account number, and gives the
credential for that account to the principal who registered it. A later
process can continue an account only if its principal holds the
credential, and only one process can run under an account at a time
(Listing~\ref{alg:stp-account-minted}).

\pdlistingroom{21}
\pdexcerpt[caption={The second version. The registry issues account
numbers, and a sanction is filed with the account it was earned
under.},label={alg:stp-account-minted}]{ch05-minted}

Carol is now refused, because she does not hold the credential for
Heron's account (Listing~\ref{alg:stp-account-run-namesake-minted}).

\pdlistingroom{6}
\pdrunlines[caption={The first break, repaired. The steps left out are the
first four of Listing~\ref{alg:stp-account-run-handover}, with an account
number where that run has a
name.},label={alg:stp-account-run-namesake-minted}]{namesake-minted}{1,10-14}{1 and 10\textendash 14}

The sanction is a different matter. \texttt{check} files a sanction with
the account it was earned under, and \texttt{register} opens a fresh
account for anyone who asks. So Alice registers the same worker again
(Listing~\ref{alg:stp-account-run-slate-minted}).

\pdlistingroom{7}
\pdrunlines[caption={The third break. The credential is authentic, and the
account is new. The steps left out are the first three of
Listing~\ref{alg:stp-account-run-slate}, with an account number where
that run has a
name.},label={alg:stp-account-run-slate-minted}]{clean-slate-minted}{1,10-15}{1 and 10\textendash 15}

This is the case Theorem~\ref{thm:necessity} describes. The account
number cannot be forged, and that does not help. A sanction filed with
one identity is lost whenever a fresh identity is free. As the chapter
puts it, an authentic credential can still be disposable.

\subsubsection{Fix: record the sanction against the principal}

Both registries file everything about a worker in one place, and that is
the mistake. \pdkernelname{} files the two kinds of entry in different
places. A pass stays with the worker's account. A sanction is recorded
against the principal, so it counts on every worker that principal runs.
The scope note under Theorem~\ref{thm:necessity} calls this a
principal-wide sanction.

The decision is made in the decision core
(\weld{prereq:sec:prereq-build}{decision core}). The kernel stores the
registry and hands it to the core with each request. The registry lists
every account with its principal, and every sanction. It also records
which account each process was started under. Nobody is asked for a
name: the kernel numbers the accounts and the processes itself, and
tells the core which process or which principal is asking.

Bob's question goes to \texttt{show}
(Listing~\ref{alg:stp-account-show}). It uses \texttt{account\_of} to
find the account the process was started under, and returns that
account's passes together with the sanctions that \texttt{count} finds.
\texttt{record\_check} decides where the result of each check goes
(Listing~\ref{alg:stp-account-check}). A pass is added to the account's
own entry. A sanction is added to the registry's list with the account's
principal on it.

\pdlistingroom{16}
\pdexcerpt[caption={The record reported for a process.},label={alg:stp-account-show}]{ch05-show}

\pdlistingroom{19}
\pdexcerpt[caption={Recording the result of a check. A pass goes on the
account, and a sanction is recorded against the
principal.},label={alg:stp-account-check}]{ch05-check}

Two guards then decide what \texttt{show} reports
(Listing~\ref{alg:stp-account-show-guards}).

\pdlistingroom{19}
\pdexcerpt[caption={The guards on what is reported. A line marked
\texttt{MUTANT} is compiled only into a test build of the core that has
that guard removed. The verification below runs each of those
builds.},label={alg:stp-account-show-guards}]{ch05-show-guards}

\texttt{count} counts every sanction recorded against the asking
process's principal, whichever worker earned it. \texttt{is\_running}
answers only the process now running under an account, and refuses any
process that ran under it earlier. Alice resets Drake again
(Listing~\ref{alg:stp-account-run-slate-core}).

\pdlistingroom{7}
\pdrunlines[caption={The second and third breaks, at the kernel's
registry. The same three steps are left
out.},label={alg:stp-account-run-slate-core}]{clean-slate-core}{1,10-15}{1 and 10\textendash 15}

The new account has no passes, and it still shows the sanction. So does
every other worker Alice runs. That is the cost of a principal-wide
sanction, and the boundary below returns to it.

\texttt{succeed} starts a new process under an existing account
(Listing~\ref{alg:stp-account-succeed}), and it has two guards of its
own (Listing~\ref{alg:stp-account-succeed-guards}).

\pdlistingroom{20}
\pdexcerpt[caption={Continuing an account with a new
process.},label={alg:stp-account-succeed}]{ch05-succeed}

\texttt{owns} lets a principal continue only the accounts it registered,
which is what refuses Carol
(Listing~\ref{alg:stp-account-run-namesake-core}). \texttt{busy} refuses
a second process while one is still running
(Listing~\ref{alg:stp-account-run-double}).

\pdlistingroom{15}
\pdexcerpt[caption={The guards on who may continue an account, and
when.},label={alg:stp-account-succeed-guards}]{ch05-succeed-guards}

\pdlistingroom{6}
\pdrunlines[caption={The first break, at the kernel's registry. The same
four steps are left
out.},label={alg:stp-account-run-namesake-core}]{namesake-core}{1,10-14}{1 and 10\textendash 14}

\pdlistingroom{12}
\pdrun[caption={One process at a time.},label={alg:stp-account-run-double}]{double-core}

A process that has been replaced may not know it, like the worker whose
lease has run out (\weld{swk:sec:swk-build}{lease}). \texttt{is\_running}
refuses it (Listing~\ref{alg:stp-account-run-replaced}).

\pdlistingroom{7}
\pdrunlines[caption={A replaced process asks for its record. Process 1
took over the account in the steps left
out.},label={alg:stp-account-run-replaced}]{replaced-core}{1,11-16}{1 and 11\textendash 16}

One more guard is not printed. When an exit is recorded, the account is
cleared only if the process that exited is the one running under it. A
late or repeated record of an old exit therefore cannot clear an account
that a newer process is using.

\keyidea[File a sanction where a fresh identity cannot escape it]{A name
is something a process says, and a credential is something it holds.
Neither tells you who answers for what the process does. A sanction
survives a reset only if it is recorded against a party that the reset
does not change. Here that party is the principal, and the worker's
passes stay with the worker's own account.}

\Needspace{12\baselineskip}
\subsubsection{Verify: every short sequence of events}

\paragraph{The chapter's table.}
Table~\ref{tab:necessity-numbers} works out by arithmetic what a reset
wins back, for records that score from 50 to 90. The kernel reproduces
the table from runs (Listing~\ref{alg:stp-account-resets}). In each run
a worker earns some passes and then a sanction: no pass for the table's
first row, and 4 for its last. Its principal then either keeps the
account or resets, and the registry is asked for the record that each
choice leaves.

% Brief (grammar: two lines on one pair of axes, with the region between
%   them filled; atlas III/fig:stp-resets).
% Reader question: after a sanction, how much of it does a reset win back,
% and above what score does a reset win back nothing?
% Claim: where a fresh identity is free, the principal ends up with the
% newcomer's 50 whatever it lost, so a reset wins back all 30 points of
% the sanction at a score of 50 and none of them from 80 up. Under the
% kernel's registry the principal ends 30 points down at every score.
% Evidence: seven records, each run at each registry [internal, harbor
% resets --each; session resets-each.txt, plotted as printed, and
% tests/harbor-research/test_stp_account_numbers.py checks the numbers in
% the caption and the prose against it].
% Must distinguish: a registry where a fresh identity is free from the
% kernel's; the score at which the two lines meet.
% Grammar chosen: the score before the sanction on the horizontal axis,
% and the best score available after it on the vertical axis, so the gap
% between the lines is what a reset wins back.
% Grammar rejected: bars for the rows of the Book's table, which show a
% gap at each row and not where the gap closes; a table, which is the
% listing beside it.
% Five-second test: where do the two lines meet?
\begin{figure}[H]
\centering
\pdfigurehue{pdviolet}%
\begin{tikzpicture}[pd figure]
  % The table is a recorded run, read as it was printed: the first line is
  % the command, and every remark starts with '#'.
  \pgfplotstableread[skip first n=1]{excerpts/resets-each.txt}\pdresets
  \begin{axis}[pd axis,width=8.2cm,height=4.5cm,
      xmin=44,xmax=116,ymin=8,ymax=92,
      xtick={50,60,70,80,90,100,110},ytick={20,50,80},
      xlabel={score before the sanction},
      ylabel={best score after it}]
    \addplot[name path=free,pd breach series] table[x=before,y=named]{\pdresets};
    \addplot[name path=kept,pd series] table[x=before,y=core]{\pdresets};
    \addplot[pd breach fill,draw=none] fill between[of=free and kept];
    \node[pd breach label,anchor=south west,align=left] at (axis cs:50.5,52) {a fresh identity\\is free};
    \node[pd label,anchor=north west] at (axis cs:86,54) {the kernel's registry};
    \node[pd label,text=pdinkmuted,anchor=west,align=left] at (axis cs:50.5,42) {won back\\by a reset};
    % Where the two lines meet: the newcomer's 50 plus the sanction's 30.
    \draw[pd guide] (axis cs:80,8) -- (axis cs:80,50);
    \node[pd note,anchor=south west] at (axis cs:80.8,9) {they meet at $50+30$};
  \end{axis}
\end{tikzpicture}
\caption{\textbf{A reset wins back the whole sanction at a score of 50, and
none of it from 80 up.} Each point is a run. Under the kernel's registry
a principal ends 30 points down whatever it does.
[internal, \texttt{harbor resets -{}-each}]}
\label{fig:stp-resets}
\end{figure}


\pdlistingroom{22}
\pdrun[caption={The chapter's table, reproduced from runs at each
registry.},label={alg:stp-account-resets}]{resets}

The registry reports counts and the table is in scores, so the run
converts one to the other. A newcomer starts at 50 and a sanction costs
30, as in the table. The table gives no value for a pass, so each adds 10
here. Under the first two registries the column \texttt{best} matches
the table's last column. A reset is worth the newcomer's 50, so it wins
back all 30 points of the sanction at a score of 50, and 10 of them at
70. Under the kernel's registry a reset is worth 20 whatever the worker
had, and wins back nothing.

Figure~\ref{fig:stp-resets} plots the same runs for seven records. Its
two lines meet at 80, which is $r_0+\sigma$, the score at which
Example~\ref{ex:newcomer-floor} finds the actor indifferent.

\paragraph{Every sequence of events.}
Those runs are a handful of cases. The kernel's exhaustive test goes
through every sequence of events for two principals, within small bounds
(Listing~\ref{alg:stp-account-walk}). The principals start up to 3
processes between them, each one a new worker or a continuation of any
existing account, their own or the other's. The kernel records up to 2
exits, possibly the same exit twice, and checks up to 2 deliveries.

\pdlistingroom{22}
\pdrun[caption={Every sequence of events for two principals, at each
registry.},label={alg:stp-account-walk}]{accounts}

After every event the test asks the registry for the record of every
process, and compares each answer with a small reference model that
keeps its own list of who started what.

The registry of names fails in every column. In 2,452 of its 11,545
states a running process shows fewer sanctions than its principal has
earned, which is the whitewashing. The registry that issues account
numbers has zeros in the other three columns, and it is still whitewashed
in 1,042 of 3,405 states. The kernel's registry has zeros in every column
across 3,717 states. It also reported every running process with the
passes it had earned, and never refused a principal its own idle account.

\paragraph{Any registry.}
That test only reaches the registries its own events can build. Kani
starts from any registry at all (\weld{anchor:sec:correct}{Kani}). It
compiles the core's real \texttt{record\_check} and \texttt{show}, and
has a solver consider every registry of a fixed size at once
(Listing~\ref{alg:stp-account-kani}).

\pdlistingroom{19}
\pdexcerpt[caption={A Kani harness on a sanction. The registry is any
registry with three accounts and four processes, with three sanctions
on it.},label={alg:stp-account-kani}]{ch05-kani}

In the harness, \texttt{any\_registry} is any registry the rules could
have produced. The solver picks the principal of each account and which
process is running under it. Any process then fails a check, and the
harness compares the record of any running process before and after. The
first assertion says that the process shows one more sanction if the
sanction's principal is its own, and the same number if not. The second
says that its passes have not changed. The 2 \texttt{cover} lines make
the solver produce both cases: a process of the same principal under a
different account, and a process of a different principal.

Nothing in these rules is a string, so Kani checks exactly the code that
runs. That was not true in the chapter before this one, where a hash had
to be shrunk for the solver (\weld{ls:sec:ls-line}{one byte}). One more
harness asserts nothing. It calls every decision on a registry that the
rules could \emph{not} have produced. Kani fails a harness if any path
through it can panic, so a pass means the core returns an answer
whatever registry it is handed
(Listing~\ref{alg:stp-account-kani-run}).

\pdlistingroom{17}
\pdrunlines[caption={Kani~0.68 on the core's rules for the
registry.},label={alg:stp-account-kani-run}]{kani-persons}{1-17}{1\textendash 17}

The fifth result is a control. Its harness counts sanctions per account,
as the second registry does, and asserts that a fresh account would
still show its principal's sanctions. That assertion is false, and Kani
reports the failure, which is what a control is for.

\paragraph{Each guard removed.}
The build has five guards. The core can be compiled with any one of them
removed, and the checks should then fail. The harnesses are run on each
of those five builds (Listing~\ref{alg:stp-account-kani-guards}), and so
is the exhaustive test (Listing~\ref{alg:stp-account-guards-walk}).

\pdlistingroom{13}
\pdrunlines[caption={The harnesses again, on five builds of the core with
one guard removed from
each.},label={alg:stp-account-kani-guards}]{kani-persons}{1,18-29}{1 and 18\textendash 29}

\pdlistingroom{20}
\pdrun[caption={The exhaustive test again, on the same five
builds.},label={alg:stp-account-guards-walk}]{accounts-guards}

Compare the rows with the breaks. A core that counts a sanction only on
the account that earned it is the second registry again, and 1,262 of its
states are whitewashed. A core that clears an account on any exit ends up
with two processes running under one account in 20 states, each of them
reached by recording the same exit twice
(Exercise~\ref{ex:stp-account-late-exit}).

Both tools catch all five. In each of the two builds before this one, one
tool missed a guard: the exhaustive test in one
(\weld{sealed:sec:sealed-gate}{two-room check}), and Kani in the other.
Here the rules are small enough for both tools to cover.

\pdscope{The exhaustive test covers two principals and 3 processes,
within the bounds printed in Listing~\ref{alg:stp-account-walk}. It
counts states and says nothing about how likely a state is. Each harness
covers one decision of the core on any registry with three accounts and
four processes.}

\begin{pdboundary}
\setlength{\parskip}{4pt plus 1pt}
A principal is only a number, the \texttt{u32} that \texttt{succeed} is
handed. The registry records a sanction against that number and never
checks who is using it. If Alice comes back under a different principal's
number, she is shown 0 sanctions. That is the same whitewashing, moved up
from the worker to the operator. Definition~\ref{def:principal} requires
a durable party, and something outside this build has to check that an
entrant is one. Inside one harbor that is the job of the cards in
\pdchapref{anchor}{The Anchor Protocol}. Across operators it is unsolved,
and \S\ref{sec:keystone} says what is missing
(Exercise~\ref{ex:stp-account-principal}).

One worker's sanction shows on every worker of the same principal.
Heron's record now carries Drake's sanction.
Definition~\ref{def:sanction-respecting} requires exactly that. It also
means that a customer comparing Alice's workers can tell them apart only
by their passes (Exercise~\ref{ex:stp-account-siblings}).

A fork is refused. The no-mint rule above lets one process have several
successors and divides its record among them. This build turns away a
second process under an account, and the no-mint rule is the chapter's
next build.

The core reports a record and computes no score. The scores in
Listing~\ref{alg:stp-account-resets} are worked out from the record
afterwards, and this build admits or refuses nothing because of a score.
\S\ref{sec:substrate} gives the reason.

A sanction is never withdrawn here. The build trusts the kernel's check,
and \S\ref{sec:revoke} deals with a recorded outcome that later proves
false.
\end{pdboundary}
